%%%%%%%%%%%%%%%%%%%%%%%%%%%%%%%%%%%%%%%%%%%%%%%%%%%%%%%%%%%%%%%%%%%%%%%%%%%%%%%
\newcommand{\faseAtual}{Projeto} % [Projeto|Final]
\newcommand{\tipo}{Conclusão de Curso} % [Iniciação Científica|Conclusão de Curso|Mestrado|Doutorado]
\newcommand{\Curso}{Ciência da Computação} % Aterar para o seu curso se necessário
\newcommand{\ano}{2026} % Alterar para o ano atual

% Dados do Orientador
% Titulação: Dr | Me | (vazio). Gênero: vazio para orientador, "a" para orientadora.
\newcommand{\orientadorNome}{Fagner de Assis}
\newcommand{\orientadorSobrenome}{Moura Pimentel}
\newcommand{\orientadorTitulo}{Dr}
\newcommand{\orientadorGenero}{}
% Coorientador (opcional: deixe o nome vazio se não houver).
\newcommand{\coorientadorNome}{}
\newcommand{\coorientadorSobrenome}{}
\newcommand{\coorientadorTitulo}{}
\newcommand{\coorientadorGenero}{}
\newcommand{\departamento}{Ciência da Computação}

% Autores (alunos): um \autor{Nome}{Sobrenome} por autor, na ordem da capa e da ficha.
% O primeiro é o autor principal da ficha catalográfica.
\newcommand{\autores}{
    \autor{Nome}{Sobrenome do Autor 1} % Alterar
    \autor{Nome}{Sobrenome do Autor 2} % Alterar, adicionar ou remover
}

% Dados do Projeto
\newcommand{\titulo}{Título do Trabalho \todo{alterar}}
\newcommand{\sub}{Subtítulo do Trabalho \todo{alterar}}

% Banca (além do orientador): um \membroBanca{Nome} por membro, na ordem das assinaturas da ata.
\newcommand{\banca}{
    \membroBanca{Prof. Dr. Fulano} % Alterar
    \membroBanca{Prof. Dr. Beltrano} % Alterar, adicionar ou remover
}

% Palavras-chave (de 3 a 5 termos, separados por vírgula).
% Usadas no resumo (Texto/resumo.tex) e na ficha catalográfica (ficha.tex).
\newcommand{\listaPalavrasChave}{Palavra-chave um, Palavra-chave dois, Palavra-chave três, Palavra-chave quatro} % Alterar


% Ficha catalográfica (não é usado em IC)
% Os campos seguem o formulário da biblioteca: https://ficha.fei.edu.br/ficha/
% Título e subtítulo vêm de \titulo e \sub acima (deixe \sub vazio se não houver).
\newcommand{\fichaCidade}{São Bernardo do Campo}
\newcommand{\fichaContagem}{f} % [f = folhas|p = páginas]
\newcommand{\fichaNumero}{20} % Alterar: número de folhas ou páginas
\newcommand{\fichaIlustracoes}{sim} % [sim|vazio]: possui ilustrações
\newcommand{\fichaFonte}{Times} % [Times|Arial]
\newcommand{\fichaTamanho}{10} % [9|10]
